\section{Conclusion}\label{section:conclusion}

While concurrent logic programming and its powerful programming techniques have been known for four decades~\cite{shapiro1989family}, adoption has been hampered by their inaccessibility to the average programmer. Concurrent logic programming requires a substantial shift from sequential procedural or functional thinking to dataflow programming via partial bindings, with concurrent operational semantics in which communication occurs through binding paired logical variables which may contain further logic variables.  Our experience to date is that AI is highly effective and productive at programming from mathematical specifications when the target is a typed high-level language such as GLP, where types constrain the space of legal programs and provide a checkable specification at the human-AI interface~\cite{shapiro2026types}. Moreover, our preliminary experience suggests that the abstract nature of concurrent logic programming renders it a more powerful and productive language than mainstream languages for collaborative human-AI abstract specifications-based program development.

Constitutional governance of grassroots communities and federations builds on Constitutional Consensus~\cite{keidar2025constitutional,lewis2025morpheus} and digital social contracts~\cite{cardelli2020digital,shapiro2022foundations}, both realisable in GLP.  The GLP runtime and example programs are open-source at \url{https://github.com/EShapiro2/GLP}.
